\section{Experimental Setup}
\label{sec:experimental_setup}

This section describes the dataset, execution environment, evaluation metrics, and experimental configurations used in our study.

% -----------------------------------------------------------------------------
\subsection{Dataset}
\label{subsec:dataset}

We evaluate all architectures on the HumanEval benchmark~\cite{chen2021humaneval}, a widely-used dataset for assessing code generation capabilities of large language models. HumanEval consists of 164 hand-written Python programming problems, each comprising a function signature with docstring, a set of assertion-based tests, and a reference solution (not used during generation).

Each task is structured as follows: the \texttt{prompt} field contains the function signature and docstring describing the expected behavior; the \texttt{test} field contains Python assertions that validate the implementation; and the \texttt{entry\_point} field specifies the function name to be implemented. This structure enables automated evaluation through assertion-based testing, where generated code is concatenated with test assertions and executed in a sandboxed subprocess with a timeout of 5 seconds.

HumanEval was chosen for several reasons: it provides a controlled benchmark with diverse problem types ranging from simple string manipulation to algorithmic challenges; it uses assertion-based testing which yields unambiguous pass/fail outcomes; and it is widely adopted in the literature, facilitating comparison with prior work.

% -----------------------------------------------------------------------------
\subsection{Execution Environment}
\label{subsec:environment}

All experiments were conducted on Kaggle Notebooks, a cloud-based computational platform providing reproducible execution environments. This choice ensures that experiments can be replicated without requiring local GPU resources.

Language models were accessed through the HuggingFace Inference API, which provides serverless endpoints for model inference. This approach offers several advantages: it eliminates the need for local model deployment, provides consistent inference performance, and enables access to models of varying sizes (1.5B to 32B parameters) without hardware constraints.

% -----------------------------------------------------------------------------
\subsection{Evaluation Metrics}
\label{subsec:metrics}

We employ four categories of metrics to comprehensively evaluate the architectures.

\paragraph{Primary Metrics.} The primary measure of functional correctness is the \textbf{Pass Rate}, defined as the percentage of tasks where all test assertions pass. We also report \textbf{Pass@1}, which measures correctness on the first generation attempt (for Architecture A this equals the pass rate, while for multi-agent architectures it reflects success before any retry or escalation). To quantify uncertainty, we compute \textbf{95\% confidence intervals} using the Wilson score interval for binomial proportions.

\paragraph{Cost Metrics.} To assess computational efficiency, we measure \textbf{Total Tokens} (sum of input and output tokens per task), \textbf{API Calls} (number of LLM invocations per task), and \textbf{Execution Time} (end-to-end wall-clock time per task in seconds). These metrics enable cost-benefit analysis of multi-agent overhead.

\paragraph{Adaptive Metrics.} For architectures employing the multi-agent pipeline (B, C, C1, C2), we track additional metrics specific to the adaptive mechanisms. \textbf{Retry Count} measures the number of Developer$\rightarrow$Tester$\rightarrow$Reviewer loops executed until success or final failure. \textbf{Escalation Count} tracks the number of tier transitions (S$\rightarrow$M or M$\rightarrow$L) triggered by test failures. \textbf{Story Point Accuracy} evaluates the correlation between the Planner's difficulty estimate and actual task outcome. \textbf{Tier Distribution} reports the percentage of tasks routed to each developer tier (S/M/L).

\paragraph{Static Code Quality Metrics.} Beyond functional correctness, we assess the quality of generated code using static analysis metrics computed via the Radon library.

\textbf{Cyclomatic Complexity (CC)}, introduced by McCabe~\cite{mccabe1976complexity}, measures the number of linearly independent paths through a program's control flow graph. It is computed as:
\begin{equation}
    CC = E - N + 2P
\end{equation}
where $E$ is the number of edges in the control flow graph, $N$ is the number of nodes, and $P$ is the number of connected components (typically 1 for a single function). Each control flow statement (\texttt{if}, \texttt{elif}, \texttt{for}, \texttt{while}, \texttt{except}, \texttt{and}, \texttt{or}) increments the complexity by one. Following Ebert et al.~\cite{ebert2016cyclomatic}, we interpret CC values as: 1--5 (low risk, simple code), 6--10 (moderate risk), 11--20 (high risk, complex code), and 21+ (very high risk, untestable).

\textbf{Maintainability Index (MI)} is a composite metric estimating how maintainable source code is~\cite{microsoft2024maintainability}. It is computed as:
\begin{equation}
    \resizebox{0.89\columnwidth}{!}{%
    $MI = \max\left(0, \frac{100 \times (171 - 5.2 \ln(HV) - 0.23 \times CC - 16.2 \ln(LOC))}{171}\right)$%
    }
\end{equation}
where $HV$ is the Halstead Volume (based on operator and operand counts), $CC$ is Cyclomatic Complexity, and $LOC$ is lines of code. The index ranges from 0 to 100, where higher values indicate better maintainability: 0--19 (low), 20--39 (moderate), 40--64 (good), 65--84 (high), and 85--100 (excellent).

These metrics are computed for all generated solutions, enabling analysis of whether multi-agent pipelines produce more maintainable code compared to single-agent approaches.

% -----------------------------------------------------------------------------
\subsection{Experimental Configuration}
\label{subsec:configuration}

We evaluate eight architectural configurations, as summarized in Table~\ref{tab:model_config}. Architectures A and A-PR represent single-agent baselines with and without prompt repetition. Architectures B and B-PR employ the full multi-agent pipeline with a single model for all roles. Architectures C and C-PR use specialized models with adaptive routing based on story points. Architecture C1 always uses the largest developer model regardless of difficulty estimate, serving as an upper-bound reference. Architecture C2 is an ablation study that removes the smallest developer tier.

Each configuration was evaluated on all 164 HumanEval tasks. For multi-agent architectures, a maximum of two escalations was permitted before terminating the pipeline. All prompt repetition variants (A-PR, B-PR, C-PR) apply the repetition technique to user prompts only, leaving system prompts unchanged.
